\bookchapter{Приложение F: Тема складных экранов}

\chapterintro{Складные экраны превратились из «технологии будущего» в зрелую категорию «можно купить сейчас». Но они сильно отличаются от обычных моноблоков по форм-фактору, привычкам использования и уходу. Стоит ли их покупать, как выбирать и как пользоваться — это приложение поможет разобраться.}

\booksection{До какой степени развились складные экраны?}

Сначала вывод:\textbf{складные экраны уже не «игрушка для ранних энтузиастов», но и не «обязательная покупка для всех» как основной форм-фактор}. Они находятся где-то между нишевыми гиковскими продуктами и массовыми потребительскими товарами и подходят определённой аудитории.

Сейчас основная масса делится на два типа:

\begin{itemize}
  \item \textbf{Горизонтальное складывание (книжный тип)}: в сложенном виде — размер обычного телефона, в разложенном — почти квадратный мини-планшет. Представители: серия Huawei Mate X, серия Samsung Galaxy Z Fold, серия Xiaomi MIX Fold, серия vivo X Fold. Основная идея: «один телефон = телефон + планшет».
  \item \textbf{Вертикальное складывание (раскладушка)}: в сложенном виде похож на пудреницу, в разложенном — обычный телефон нормального размера. Представители: серия Huawei Pocket, серия Samsung Galaxy Z Flip, OPPO Find N Flip, Xiaomi MIX Flip. Основная идея: «компактность + быстрый просмотр на внешнем экране».
\end{itemize}

Ключевые достижения 2024–2026 годов: ресурс петель обычно заявлен на 500 000 складываний (при 100 складываний в день хватит на 13 лет), UTG сверхтонкое стекло пришло на смену ранним пластиковым плёнкам, складка видна глазу, но на ощупь уже едва заметна, цена опустилась с десятитысячного диапазона до 4000–5000 юаней.

\begin{tcbinfo}{Краткий вывод}
Если вам полностью хватает обычного моноблока и вы не испытываете «боли от недостаточно большого экрана» — не покупайте складной. Это устройство «дополняющее», а не «спасательное».

\end{tcbinfo}

\booksection{Кому подходят складные экраны?}

По реальным отзывам пользователей, этим группам складные экраны с большой вероятностью понравятся:

\begin{enumerate}
  \item \textbf{Активные читатели / те, кто работает с документами}: часто просматривают PDF на телефоне, аннотируют документы, читают длинные тексты. В разложенном виде внутренний экран почти 7.5 дюйма — читать A4 PDF без масштабирования, опыт кардинально меняется.
  \item \textbf{Деловые пользователи, часто в командировках / на встречах}: не хотят брать планшет, но телефон кажется слишком маленьким. На горизонтальном складном экране можно смотреть презентации, подписывать договоры, сравнивать письма на двух экранах — заметно эффективнее обычного телефона.
  \item \textbf{Любители творчества в «режиме полустоя»}: сложите на 90 градусов и поставьте на стол — верхний экран для кадра, нижний для спуска (селфи/таймлапс), верхний для видео, нижний для заметок, верхний для стрима, нижний для комментариев — такой форм-фактор недоступен на обычном телефоне.
  \item \textbf{Женщины, которым важна и компактность, и большой экран}: вертикальный складной в сложенном виде помещается в маленькую сумочку, в разложенном — обычный телефон, а на внешнем экране можно пользоваться зеркальцем и смотреть уведомления.
\end{enumerate}

\begin{tcbwarn}{Этим группам лучше не покупать}
\begin{itemize}
  \item Тем, кто использует телефон только для коротких видео, переписки в WeChat и звонков — складной экран не даст прироста, а будет лишь тяжелее и толще.
  \item С небольшой ладонью или тендинитом запястья — горизонтальный складной в разложенном виде обычно весит 230 г+, держать одной рукой утомительно.
  \item При ограниченном бюджете, если телефон планируется использовать 4+ года — у складных экранов остаточная стоимость ниже, чем у моноблоков, а стоимость ремонта выше.
\end{itemize}

\end{tcbwarn}

\booksection{На что смотреть при выборе: 6 показателей}

Не дайте себя ослепить словом «складной», при выборе обращайте внимание на:

\begin{itemize}
  \item \textbf{Тип петли}: капельная петля (складка неглубокая, ресурс большой, сейчас основная) превосходит U-образную петлю (складка глубокая, раннее решение). Спросите продавца или проверьте обзоры, чтобы уточнить тип.
  \item \textbf{Материал внутреннего экрана}: обязательно UTG сверхтонкое стекло, а не ранняя пластиковая плёнка. UTG на ощупь близок к стеклу, по пластиковой плёнке ноготь оставляет вдавленность и следы.
  \item \textbf{Проявление складки}: сходите в офлайн-магазин и сложите/разложите сами. К 2026 году приемлемый уровень: при включённом экране спереди почти незаметно, под углом видна неглубокая бороздка, при проводе пальцем — лёгкая впадина, но не режет руку.
  \item \textbf{Размер внешнего экрана}: у горизонтальных складных чем больше внешний экран, тем лучше (желательно 6.4 дюйма и выше), чтобы в сложенном виде пользоваться как обычным телефоном; у вертикальных складных чем больше внешний экран, тем практичнее (4 дюйма и выше, чтобы отвечать на сообщения, смотреть навигацию).
  \item \textbf{Вес и толщина}: для горизонтальных складных в разложенном виде 230–260 г — норма, менее 230 г — отлично; толщина в сложенном виде 12–14 мм, менее 12 мм — тонко. Для вертикальных складных важно, чтобы в сложенном виде легко помещался в карман рубашки.
  \item \textbf{Сервис и гарантия}: в первую очередь уточняйте стоимость ремонта при разбитии внутреннего экрана. Некоторые бренды в первый год бесплатно меняют внутренний экран один раз, у других замена обойдётся в 2000–4000 юаней. Это самая большая скрытая стоимость складных телефонов.
\end{itemize}

\begin{tcbpractice}{Чек-лист для офлайн-тестирования}
\begin{enumerate}
  \item Самостоятельно сложите/разложите 10 раз, оцените плавность хода петли и наличие посторонних звуков.
  \item В разложенном виде посмотрите на складку под углом к свету, затем спереди при включённом экране.
  \item Удерживайте устройство в разложенном виде одной рукой 30 секунд, оцените, приемлем ли вес.
  \item Попробуйте режим полустоя: сложите на 90 градусов, положите на стол, попробуйте сфотографировать и посмотреть видео.
  \item Уточните стоимость отдельного ремонта внутреннего экрана и условия гарантии.
\end{enumerate}

\end{tcbpractice}

\booksection{Приёмы использования, доступные только на складных экранах}

После покупки эти сценарии недоступны на обычных телефонах:

\booksubsection{Разделение экрана и многозадачность}

На горизонтальном складном экране в разложенном виде большинство систем поддерживает верхнее/нижнее или левое/правое разделение. Типовые сценарии: верхний экран — видео онлайн-курса, нижний — заметки; слева WeChat, справа документ; сверху карта-навигатор, снизу музыка. Это намного практичнее, чем на обычном телефоне, потому что экран достаточно большой, и каждая половина не тесная.

\booksubsection{Режим полустоя (складывание на 90 градусов)}

Это «киллер-фича» складных экранов:

\begin{itemize}
  \item \textbf{Селфи / групповое фото}: поставьте на стол, сложите на 90 градусов, верхний экран для кадра, нижний — для спуска, не нужно просить кого-то сфотографировать.
  \item \textbf{Таймлапс / ночная съёмка}: так же на столе для стабильной съёмки, без штатива.
  \item \textbf{Видеозвонок}: освобождает руки, ракурс, под которым вас видит собеседник, тоже более естественный.
  \item \textbf{Просмотр видео}: поставьте на стол как маленький телевизор, особенно удобно на откидных столиках в самолёте / поезде.
  \item \textbf{Предпросмотр на внешнем экране + съёмка на внутреннем}: при съёмке себя основной камерой внешний экран в реальном времени показывает кадр — точнее компонуйте.
\end{itemize}

\booksubsection{Быстрый просмотр на внешнем экране}

У вертикального складного телефона в сложенном виде на внешнем экране можно: смотреть время / уведомления / погоду, отвечать в WeChat (преобразование голоса в текст или быстрый ответ), прикладывать транспортную карту, смотреть навигацию, использовать как зеркальце (вызывая фронтальную камеру). Не нужно каждый раз раскладывать — экономит заряд и время.

\booksubsection{Режим компьютера}

Некоторые горизонтальные складные (Huawei, Samsung) после раскрытия могут переключаться в «режим компьютера»: при подключении внешнего монитора, клавиатуры и мыши работает как простенький ПК — для временной правки документов и ответов на письма достаточно. Это не основной инструмент производительности, но выручит в экстренной ситуации.

\booksection{Уход: как продлить жизнь складному экрану}

Самые хрупкие элементы складного телефона — внутренний экран и петля. Эти привычки значительно продлевают срок службы:

\begin{itemize}
  \item \textbf{Не складывайте/раскладывайте при сильном морозе}: при температуре ниже 0 °C UTG-стекло становится хрупким, перед складыванием подержите устройство несколько минут у тела. Это пишут во всех официальных инструкциях, но мало кто их читает.
  \item \textbf{Избегайте контакта ногтей и твёрдых предметов с внутренним экраном}: хотя у внутреннего экрана есть защитный слой, он мягче внешнего; работайте подушечкой пальца, не тыкайте ногтем.
  \item \textbf{Перед складыванием убедитесь, что внутри нет посторонних предметов}: одна песчинка, зажатая при складывании, может прямо повредить внутренний экран. Перед складыванием окиньте взглядом.
  \item \textbf{Не наклеивайте обычную защитную плёнку}: для внутреннего экрана подходит только официальная плёнка или UTG-защитная плёнка; обычная закалённая плёнка влияет на складывание и сенсор. Рекомендуется наклеивать в официальном сервисе.
  \item \textbf{Не допускайте попадания пыли и жидкости в петлю}: хотя большинство моделей заявляют водозащиту IPX8, пылезащита обычно невысокая. Будьте осторожны на пляже, стройке и т. п.
  \item \textbf{Складывайте/раскладывайте равномерно}: не раскрывайте резко одной рукой; наиболее надёжно раскрывать двумя руками равномерно от центра.
\end{itemize}

\textbf{Типичные сценарии повреждения внутреннего экрана}\par

\begin{itemize}
  \item При складывании зажать провод наушников / зарядки — внутренний экран трескается прямо.
  \item В сумке вперемешку с ключами — внешний экран цел, а на внутреннем появляются царапины от ключей.
  \item Зимой зайти с мороза в тёплое помещение и сразу начать складывать/раскладывать — перепад температур вызывает изменение напряжения стекла.
  \item Через год эксплуатации решить, что «складка стала глубже», и отклеить заводскую плёнку — снятие плёнки повреждает и сам внутренний экран.
\end{itemize}

\booksection{«Подводные камни» складных экранов и ожидания}

Перед покупкой держите в уме эти моменты:

\begin{itemize}
  \item \textbf{Складка будет всегда}: технологии 2026 года не позволяют сделать «полностью без следов», можно лишь добиться «спереди при включённом экране почти незаметно». Если вы перфекционист и вас каждый день будет раздражать складка — не покупайте.
  \item \textbf{Тяжелее и толще}: горизонтальный складной в разложенном виде обычно весит 230 г+, держать одной рукой утомительнее, чем обычный телефон; в сложенном виде толщина 12–14 мм, в кармане брюк выпирает сильнее.
  \item \textbf{Автономность обычно средняя}: из-за необходимости разместить механизм петли место под батарею ограничено, при той же цене автономность обычно уступает обычным телефонам.
  \item \textbf{Низкая остаточная стоимость}: через год цена на вторичном рынке обычно 50–60 \% от цены нового, падает быстрее, чем у обычных телефонов.
  \item \textbf{Дорогой ремонт}: замена модуля внутреннего экрана — 2000–4000 юаней, ремонт петли — 500–1500 юаней. Перед покупкой уточните условия гарантии.
  \item \textbf{Неравномерная адаптация приложений}: основные приложения уже адаптированы под большой экран, но нишевые могут просто «растягиваться», опыт посредственный.
\end{itemize}

Складной экран — это не «лучший телефон», а «телефон другой формы». Его покупают потому, что нужен именно этот уникальный форм-фактор, а не потому, что он «продвинутее».

— Философия выбора складных экранов

\booksection{Какой бренд выбрать (2026 год)}

Не приводим конкретные модели (они устареют), описываем только фирменный стиль, чтобы помочь быстро сориентироваться:

\begin{itemize}
  \item \textbf{Серия Huawei Mate X}: эталон горизонтального складывания, сильная синергия экосистемы HarmonyOS, пропорции внешнего экрана близки к обычному телефону, в сложенном виде можно использовать как основной аппарат. Цена высокая, но остаточная стоимость относительно лучшая.
  \item \textbf{Samsung Galaxy Z Fold/Flip}: первопроходец складных экранов, отлаженные петли, Z Flip — лидер по дизайну среди вертикальных. Недостаток — небольшая сеть сервисных центров в Китае.
  \item \textbf{Xiaomi MIX Fold/Flip}: маршрут «цена/качество», богатая комплектация, относительно доступная цена. Подходит тем, кто хочет попробовать, но не готов тратить много.
  \item \textbf{Серия vivo X Fold}: среди горизонтальных складных имеет пропорции внешнего экрана, наиболее близкие к обычному телефону, удобно использовать в сложенном виде как моноблок.
  \item \textbf{Серия OPPO Find N}: среди горизонтальных складных более компактный размер, можно держать одной рукой — для тех, кому большие складные слишком тяжёлые.
  \item \textbf{Серия Huawei Pocket}: среди вертикальных складных самый большой и практичный внешний экран, тщательно проработанные женские функции — зеркальце, проверка защиты от солнца и т. п.
\end{itemize}

Совет по выбору: при большом бюджете и желании экосистемы — Huawei Mate X; за самой отлаженной петлёй — Samsung Z Fold; лучшая цена/качество — Xiaomi/vivo; для вертикального складывания — Huawei Pocket и Samsung Z Flip.

\begin{tcbinfo}{Напоследок}
Технологии складных экранов развиваются с каждым годом, это приложение написано в 2026 году. Перед покупкой обязательно изучите свежие обзоры и отзывы пользователей, обращая внимание на: реальный ресурс петли, проявление складки и условия ремонта внутреннего экрана. Потрогать устройство вживую всегда важнее, чем читать обзоры.

\end{tcbinfo}
